% Options for packages loaded elsewhere
\PassOptionsToPackage{unicode}{hyperref}
\PassOptionsToPackage{hyphens}{url}
\PassOptionsToPackage{dvipsnames,svgnames,x11names}{xcolor}
\documentclass[
  french,
  11pt,
  a4paper,
]{article}
\usepackage{xcolor}
\usepackage[margin=2.3cm]{geometry}
\usepackage{amsmath,amssymb}
\setcounter{secnumdepth}{5}
\usepackage{iftex}
\ifPDFTeX
  \usepackage[T1]{fontenc}
  \usepackage[utf8]{inputenc}
  \usepackage{textcomp} % provide euro and other symbols
\else % if luatex or xetex
  \usepackage{unicode-math} % this also loads fontspec
  \defaultfontfeatures{Scale=MatchLowercase}
  \defaultfontfeatures[\rmfamily]{Ligatures=TeX,Scale=1}
\fi
\usepackage{lmodern}
\ifPDFTeX\else
  % xetex/luatex font selection
  \setmainfont[]{TeX Gyre Pagella}
  \setmathfont[]{TeX Gyre Pagella Math}
\fi
% Use upquote if available, for straight quotes in verbatim environments
\IfFileExists{upquote.sty}{\usepackage{upquote}}{}
\IfFileExists{microtype.sty}{% use microtype if available
  \usepackage[]{microtype}
  \UseMicrotypeSet[protrusion]{basicmath} % disable protrusion for tt fonts
}{}
\makeatletter
\@ifundefined{KOMAClassName}{% if non-KOMA class
  \IfFileExists{parskip.sty}{%
    \usepackage{parskip}
  }{% else
    \setlength{\parindent}{0pt}
    \setlength{\parskip}{6pt plus 2pt minus 1pt}}
}{% if KOMA class
  \KOMAoptions{parskip=half}}
\makeatother
\usepackage{longtable,booktabs,array}
\newcounter{none} % for unnumbered tables
\usepackage{calc} % for calculating minipage widths
% Correct order of tables after \paragraph or \subparagraph
\usepackage{etoolbox}
\makeatletter
\patchcmd\longtable{\par}{\if@noskipsec\mbox{}\fi\par}{}{}
\makeatother
% Allow footnotes in longtable head/foot
\IfFileExists{footnotehyper.sty}{\usepackage{footnotehyper}}{\usepackage{footnote}}
\makesavenoteenv{longtable}
\ifLuaTeX
\usepackage[bidi=basic,shorthands=off]{babel}
\else
\usepackage[bidi=default,shorthands=off]{babel}
\fi
\ifPDFTeX
\else
\babelfont{rm}[]{TeX Gyre Pagella}
\fi
\ifLuaTeX
  \usepackage{selnolig} % disable illegal ligatures
\fi
\setlength{\emergencystretch}{3em} % prevent overfull lines
\providecommand{\tightlist}{%
  \setlength{\itemsep}{0pt}\setlength{\parskip}{0pt}}
\usepackage[style=apa]{biblatex}
\addbibresource{../../02_ETAT_DE_L_ART_FINAL/references.bib}
\usepackage{etoolbox,caption}\captionsetup{font=small,labelfont=bf,labelsep=endash}\captionsetup[table]{position=top}\AtBeginEnvironment{longtable}{\small}\addto\captionsfrench{\renewcommand{\tablename}{Tableau}}
\usepackage{bookmark}
\IfFileExists{xurl.sty}{\usepackage{xurl}}{} % add URL line breaks if available
\urlstyle{same}
\hypersetup{
  pdftitle={Analyse des résultats expérimentaux --- RL-in-GA pour le problème du sac à dos},
  pdfauthor={RALAIVAO Niaiko Michaël --- Master Recherche, ENI / LIMAD / GLoRIA},
  pdflang={fr},
  colorlinks=true,
  linkcolor={blue},
  filecolor={Maroon},
  citecolor={teal},
  urlcolor={blue},
  pdfcreator={LaTeX via pandoc}}

\title{Analyse des résultats expérimentaux --- RL-in-GA pour le problème
du sac à dos}
\author{RALAIVAO Niaiko Michaël --- Master Recherche, ENI / LIMAD /
GLoRIA}
\date{Septembre 2026}

\begin{document}
\maketitle

{
\setcounter{tocdepth}{3}
\tableofcontents
}
\section{Organisation des
résultats}\label{analyse_resultats-organisation-des-ruxe9sultats}

Les résultats sont regroupés par campagne ; chaque figure porte un
numéro stable (préfixe \texttt{figNN\_}) et chaque tableau existe en
CSV, Markdown et LaTeX dans
\texttt{results/analyse/\textless{}campagne\textgreater{}/}, pour une
intégration directe dans le manuscrit.

{\def\LTcaptype{none} % do not increment counter
\begin{longtable}[]{@{}
  >{\raggedright\arraybackslash}p{(\linewidth - 6\tabcolsep) * \real{0.2500}}
  >{\raggedright\arraybackslash}p{(\linewidth - 6\tabcolsep) * \real{0.2500}}
  >{\raggedright\arraybackslash}p{(\linewidth - 6\tabcolsep) * \real{0.2500}}
  >{\raggedright\arraybackslash}p{(\linewidth - 6\tabcolsep) * \real{0.2500}}@{}}
\toprule\noalign{}
\begin{minipage}[b]{\linewidth}\raggedright
Campagne
\end{minipage} & \begin{minipage}[b]{\linewidth}\raggedright
Question
\end{minipage} & \begin{minipage}[b]{\linewidth}\raggedright
Tableaux (\texttt{results/analyse/…})
\end{minipage} & \begin{minipage}[b]{\linewidth}\raggedright
Figures (\texttt{results/figures/…})
\end{minipage} \\
\midrule\noalign{}
\endhead
\bottomrule\noalign{}
\endlastfoot
E0 --- calibration & quels réglages ? & \texttt{exp0\_pilot/} &
\texttt{fig00\_*}, \texttt{fig05\_*}, \texttt{fig06\_*} \\
E1 --- KP générés & QR1, QR2, QR3, QR5 & \texttt{e1\_kp01\_generes/} &
\texttt{fig01\_e1\_*}, \texttt{fig04\_politique\_actions} \\
E2 --- KP de référence & QR1, QR3, QR4 & \texttt{e2\_kp01\_benchmarks/}
& \texttt{fig02\_e2\_*} \\
E3 --- MKP & QR3 (transfert inter-problèmes), QR4 & \texttt{e3\_mkp/} &
\texttt{fig03\_e3\_*} \\
E5 --- passage à l'échelle & QR4 &
\texttt{e5\_scalabilite/scalability.csv} &
\texttt{fig07\_scalabilite} \\
Cas d'usage & utilité pratique &
\texttt{cas\_usage/resultats\_cas\_usage.csv} & \texttt{app\_*}
(captures) \\
\end{longtable}
}

\textbf{Lecture des indicateurs.} L'\emph{écart} est
\(100\,(f_{\text{ref}} - f)/f_{\text{ref}}\) en \% (0 = optimum). Les
écarts sont faibles en valeur absolue (quelques centièmes de \%) parce
que l'AG, grâce à la réparation gloutonne, est déjà un bon solveur ; ce
sont les \textbf{différences relatives} entre méthodes et leur
\textbf{significativité} qui importent. Les comparaisons suivent les
recommandations de \textcite{demsar_2006} et \textcite{derrac_2011} :
test de Friedman suivi du test post hoc de Nemenyi pour le classement
global, test de Wilcoxon apparié avec correction de Holm pour les
comparaisons deux à deux. \emph{V/E/D} = nombre d'instances où la
méthode cible gagne / fait jeu égal / perd (Mann--Whitney, α = 0,05) ;
\(A_{12} > 0{,}5\) signifie que la cible obtient plus souvent un écart
plus faible (taille d'effet de \textcite{vargha_delaney_2000}).

\section{E0 --- Calibration (instances disjointes des
tests)}\label{analyse_resultats-e0-calibration-instances-disjointes-des-tests}

Réalisée sur 6 instances (SC, ASC, UC ; n = 500 ; graines 900+), selon
la démarche de réglage des paramètres des algorithmes évolutionnaires
décrite par \textcite{eiben_1999} :

\begin{enumerate}
\def\labelenumi{\arabic{enumi}.}
\tightlist
\item
  \textbf{Tous les contrôleurs adaptatifs} (aléatoire, UCB1, QL en
  ligne) divisent par ≈ 4 l'écart de l'AG standard sur SC/ASC, mais le
  \textbf{QL en ligne ne dépasse pas la sélection aléatoire} (écart
  moyen 0,0182 \% contre 0,0158 \%).
\item
  Parmi les \textbf{9 configurations fixes}, \textbf{UX + BF3} est de
  loin la meilleure (SC : 0,0000 \% ; ASC : 0,0073 \%) et 1P + BF1 (l'AG
  standard) la pire (SC : 0,1105 \%) --- un facteur \textgreater{} 10 :
  \textbf{le choix des opérateurs compte}, et il existe donc une marge
  pour l'apprentissage, ce qui est la prémisse de la sélection
  adaptative d'opérateurs
  \autocite{dacosta_2008,karimi-mamaghan_machine_2022}.
\item
  Choix de la \textbf{récompense} (Q-table entraînée sur UC/WC/SS puis
  testée) : l'amélioration relative seule (R1) donne la meilleure
  politique (SC 0,0000 \%, ASC 0,0149 \%) ; l'ajout du taux de succès
  des descendants (SC 0,0237 \%) ou d'une pénalité de diversité (SC
  0,0395 \%) la dégrade. → \textbf{R1 retenue}. Ce résultat est cohérent
  avec les récompenses fondées sur l'amélioration de la qualité, les
  plus fréquentes dans les revues sur l'apprentissage par renforcement
  au service des algorithmes évolutionnaires
  \autocite{song_reinforcement_2023,li_bridging_2025}.
\end{enumerate}

\section{E1 --- KP 0/1 générés (54 instances × 7 variantes × 10
exécutions)}\label{analyse_resultats-e1-kp-01-guxe9nuxe9ruxe9s-54-instances-7-variantes-10-exuxe9cutions}

Les instances sont générées selon les classes de
\textcite{pisinger_where_2005} ; les contrôleurs comparés sont l'AG
standard \autocite{holland_1975,goldberg_1989}, la sélection aléatoire,
le bandit UCB1 \autocite{auer_2002} et le Q-learning
\autocite{watkins_dayan_1992,sutton_barto_2018}, en ligne ou
pré-entraîné.

\textbf{Tableau 1 --- Écart moyen à l'optimum (\%)} (optimum prouvé par
programmation dynamique
\autocite{bellman_1957,kellerer_knapsack_book_2004}) :

{\def\LTcaptype{none} % do not increment counter
\begin{longtable}[]{@{}
  >{\raggedright\arraybackslash}p{(\linewidth - 14\tabcolsep) * \real{0.1250}}
  >{\raggedright\arraybackslash}p{(\linewidth - 14\tabcolsep) * \real{0.1250}}
  >{\raggedright\arraybackslash}p{(\linewidth - 14\tabcolsep) * \real{0.1250}}
  >{\raggedright\arraybackslash}p{(\linewidth - 14\tabcolsep) * \real{0.1250}}
  >{\raggedright\arraybackslash}p{(\linewidth - 14\tabcolsep) * \real{0.1250}}
  >{\raggedright\arraybackslash}p{(\linewidth - 14\tabcolsep) * \real{0.1250}}
  >{\raggedright\arraybackslash}p{(\linewidth - 14\tabcolsep) * \real{0.1250}}
  >{\raggedright\arraybackslash}p{(\linewidth - 14\tabcolsep) * \real{0.1250}}@{}}
\toprule\noalign{}
\begin{minipage}[b]{\linewidth}\raggedright
Classe
\end{minipage} & \begin{minipage}[b]{\linewidth}\raggedright
AG standard
\end{minipage} & \begin{minipage}[b]{\linewidth}\raggedright
Aléatoire
\end{minipage} & \begin{minipage}[b]{\linewidth}\raggedright
UCB1
\end{minipage} & \begin{minipage}[b]{\linewidth}\raggedright
QL en ligne
\end{minipage} & \begin{minipage}[b]{\linewidth}\raggedright
\textbf{QL pré-entraîné}
\end{minipage} & \begin{minipage}[b]{\linewidth}\raggedright
QL pré-entr. + φ(I)
\end{minipage} & \begin{minipage}[b]{\linewidth}\raggedright
Oracle fixe
\end{minipage} \\
\midrule\noalign{}
\endhead
\bottomrule\noalign{}
\endlastfoot
UC & 0,0048 & 0,0028 & 0,0024 & 0,0031 & \textbf{0,0019} & 0,0031 &
0,0021 \\
WC & 0,0063 & 0,0039 & 0,0036 & 0,0039 & 0,0037 & 0,0035 &
\textbf{0,0033} \\
SS & 0,0000 & 0,0000 & 0,0000 & 0,0000 & 0,0000 & 0,0000 & 0,0000 \\
ISC & 0,0053 & 0,0000 & 0,0000 & 0,0000 & 0,0000 & 0,0000 & 0,0000 \\
SC & 0,1114 & 0,0258 & 0,0275 & 0,0275 & 0,0158 & 0,0258 &
\textbf{0,0086} \\
ASC & 0,1016 & 0,0341 & 0,0321 & 0,0319 & 0,0232 & 0,0309 &
\textbf{0,0156} \\
\textbf{Moyenne} & 0,0382 & 0,0111 & 0,0109 & 0,0111 & \textbf{0,0074} &
0,0105 & 0,0049 \\
Rang moyen (Friedman) & 6,05 & 4,08 & 3,96 & 4,08 & \textbf{3,10} & 4,09
& 2,63 \\
\end{longtable}
}

Test de Friedman : χ² = 130,15, p = 1,2·10⁻²⁵ (N = 54 instances,
différence critique de Nemenyi CD = 1,23). Pour comparaison, le
\textbf{glouton} fondé sur l'ordre des efficacités
\autocite{dantzig_1957} laisse un écart moyen de 0,34 \% sur SC et ASC,
soit \textbf{plus de 10 fois} celui des variantes adaptatives de l'AG.

\textbf{Tableau 2 --- Tests appariés (Wilcoxon sur les 54 écarts moyens
par instance, correction de Holm)} :

{\def\LTcaptype{none} % do not increment counter
\begin{longtable}[]{@{}
  >{\raggedright\arraybackslash}p{(\linewidth - 8\tabcolsep) * \real{0.2000}}
  >{\raggedright\arraybackslash}p{(\linewidth - 8\tabcolsep) * \real{0.2000}}
  >{\raggedright\arraybackslash}p{(\linewidth - 8\tabcolsep) * \real{0.2000}}
  >{\raggedright\arraybackslash}p{(\linewidth - 8\tabcolsep) * \real{0.2000}}
  >{\raggedright\arraybackslash}p{(\linewidth - 8\tabcolsep) * \real{0.2000}}@{}}
\toprule\noalign{}
\begin{minipage}[b]{\linewidth}\raggedright
Comparaison
\end{minipage} & \begin{minipage}[b]{\linewidth}\raggedright
Écarts (\%)
\end{minipage} & \begin{minipage}[b]{\linewidth}\raggedright
p (Holm)
\end{minipage} & \begin{minipage}[b]{\linewidth}\raggedright
\(A_{12}\)
\end{minipage} & \begin{minipage}[b]{\linewidth}\raggedright
V/E/D
\end{minipage} \\
\midrule\noalign{}
\endhead
\bottomrule\noalign{}
\endlastfoot
QL pré-entraîné vs AG standard & 0,0074 vs 0,0382 & \textless{} 10⁻⁵ &
0,74 & 28/26/0 \\
QL pré-entraîné vs aléatoire & 0,0074 vs 0,0111 & 0,0005 & 0,59 &
9/45/0 \\
QL pré-entraîné vs UCB1 & 0,0074 vs 0,0109 & 0,003 & 0,58 & 10/44/0 \\
QL pré-entraîné vs QL en ligne & 0,0074 vs 0,0111 & 0,002 & 0,58 &
7/47/0 \\
QL pré-entraîné vs oracle fixe & 0,0074 vs 0,0049 & 0,003 & 0,42 &
0/45/9 \\
QL en ligne vs AG standard & 0,0111 vs 0,0382 & \textless{} 10⁻⁵ & 0,71
& 24/30/0 \\
QL en ligne vs aléatoire & 0,0111 vs 0,0111 & 1,00 (n.s.) & 0,51 &
0/53/1 \\
QL en ligne vs UCB1 & 0,0111 vs 0,0109 & 1,00 (n.s.) & 0,51 & 1/51/2 \\
\end{longtable}
}

\textbf{Effet de la taille} (classes SC et ASC) : l'écart de l'AG
standard croît de 0,046 \% (n = 200) à 0,167 \% (n = 1 000) ; celui du
QL pré-entraîné passe de 0,005 \% à 0,025 \%, contre 0,043 \% pour le QL
en ligne et l'aléatoire à n = 1 000. \textbf{L'avantage du
pré-entraînement augmente avec la taille} : les instances de 500 et 1
000 objets sont hors de la distribution d'entraînement (≤ 200 objets).

\textbf{Vitesse} : le QL pré-entraîné atteint son meilleur résultat plus
tôt (génération 48,8 en moyenne, contre 56,9 pour l'AG standard et 50,5
pour l'aléatoire) ; le surcoût de l'agent est négligeable (≈ 4,6 s
contre 4,0 s par exécution, à n moyen de 567).

\textbf{Méthodes exactes} : la PD résout toutes les instances d'E1 en ≤
3,4 s ; la PLNE (solveur HiGHS \autocite{huangfu_hall_2018}, appelé
depuis SciPy \autocite{virtanen_scipy_2020}) trouve l'optimum sur les 54
instances mais \textbf{ne prouve pas l'optimalité} dans la limite de 30
s pour 12 instances corrélées (SC, ISC, ASC de 500--1 000 objets),
confirmant leur difficulté pour le \emph{Branch-and-Bound}
\autocite{pisinger_where_2005,martello_toth_1990}.

\textbf{Figures associées} : \texttt{fig01\_e1\_ecart\_par\_classe}
(écarts et IC 95 \%), \texttt{fig01\_e1\_rangs\_friedman},
\texttt{fig01\_e1\_boites} (distributions),
\texttt{fig01\_e1\_convergence} (courbes et diversité),
\texttt{fig01\_e1\_succes\_echecs}, \texttt{fig01\_e1\_qualite\_temps},
\texttt{fig04\_politique\_actions}.

\subsection{Ce que montre la politique apprise
(QR5)}\label{analyse_resultats-ce-que-montre-la-politique-apprise-qr5}

La figure \texttt{fig04\_politique\_actions} donne la fréquence des 9
actions par tiers de la recherche :

\begin{itemize}
\tightlist
\item
  \textbf{QL en ligne et UCB1} : distributions \textbf{quasi uniformes}
  (≈ 11 \% par action dans toutes les phases). En 200 générations, ils
  n'identifient pas de préférence : leur comportement est en pratique
  celui de la sélection aléatoire --- ce qui \textbf{explique} l'égalité
  statistique avec GA-RAND (H2).
\item
  \textbf{QL pré-entraîné} : préférence nette pour le \textbf{croisement
  uniforme} --- UX + SWAP (40 \%) et UX + BF3 (17--30 \%) en début de
  recherche, UX + BF3 (46--54 \%) et UX + SWAP (33--39 \%) en fin de
  recherche --- soit exactement la famille de l'oracle (UX + BF3),
  \textbf{découverte sans la connaître} --- ce qui illustre l'intérêt de
  la configuration dynamique d'algorithmes apprise hors ligne
  \autocite{biedenkapp_2020,benjamins_2024} ---, sur des instances
  faciles, et appliquée aux classes difficiles. En phase intermédiaire,
  la répartition est presque uniforme : ces états ont été peu visités
  pendant l'entraînement (valeurs Q quasi nulles), ce qui explique
  l'écart résiduel avec l'oracle.
\item
  \textbf{φ(I)} n'apporte rien ici : avec 81 états au lieu de 27 pour un
  budget d'entraînement comparable, chaque état est moins visité ; la
  politique est plus diluée (écart 0,0105 \% contre 0,0074 \%). C'est la
  « malédiction de la dimension » des représentations tabulaires, qui
  motive l'approximation de fonction
  \autocite{sutton_barto_2018,mnih_2015}.
\end{itemize}

\section{\texorpdfstring{E2 --- KP 0/1 de référence :
\texttt{instances\_01\_KP} et Jooken et al.~(31 instances × 6 variantes
× 10
exécutions)}{E2 --- KP 0/1 de référence : instances\_01\_KP et Jooken et al.~(31 instances × 6 variantes × 10 exécutions)}}\label{analyse_resultats-e2-kp-01-de-ruxe9fuxe9rence-instances_01_kp-et-jooken-et-al.-31-instances-6-variantes-10-exuxe9cutions}

Ce jeu réunit des instances de Pisinger et les instances difficiles de
\textcite{jooken_2022}, conçues pour mettre en défaut les méthodes
exactes.

\textbf{Tableau 3 --- Écart moyen à l'optimum (\%)} (optimums fournis
avec les instances) :

{\def\LTcaptype{none} % do not increment counter
\begin{longtable}[]{@{}
  >{\raggedright\arraybackslash}p{(\linewidth - 12\tabcolsep) * \real{0.1429}}
  >{\raggedright\arraybackslash}p{(\linewidth - 12\tabcolsep) * \real{0.1429}}
  >{\raggedright\arraybackslash}p{(\linewidth - 12\tabcolsep) * \real{0.1429}}
  >{\raggedright\arraybackslash}p{(\linewidth - 12\tabcolsep) * \real{0.1429}}
  >{\raggedright\arraybackslash}p{(\linewidth - 12\tabcolsep) * \real{0.1429}}
  >{\raggedright\arraybackslash}p{(\linewidth - 12\tabcolsep) * \real{0.1429}}
  >{\raggedright\arraybackslash}p{(\linewidth - 12\tabcolsep) * \real{0.1429}}@{}}
\toprule\noalign{}
\begin{minipage}[b]{\linewidth}\raggedright
Classe
\end{minipage} & \begin{minipage}[b]{\linewidth}\raggedright
AG standard
\end{minipage} & \begin{minipage}[b]{\linewidth}\raggedright
Aléatoire
\end{minipage} & \begin{minipage}[b]{\linewidth}\raggedright
UCB1
\end{minipage} & \begin{minipage}[b]{\linewidth}\raggedright
QL en ligne
\end{minipage} & \begin{minipage}[b]{\linewidth}\raggedright
\textbf{QL pré-entraîné}
\end{minipage} & \begin{minipage}[b]{\linewidth}\raggedright
Oracle fixe
\end{minipage} \\
\midrule\noalign{}
\endhead
\bottomrule\noalign{}
\endlastfoot
UC (n = 100\ldots2 000) & 0,0345 & \textbf{0,0184} & 0,0221 & 0,0190 &
0,0222 & 0,0229 \\
WC & 0,1904 & 0,1161 & 0,1571 & 0,1381 & \textbf{0,0827} & 0,1144 \\
SC & 0,0232 & 0,0142 & \textbf{0,0001} & 0,0008 & \textbf{0,0001} &
0,0139 \\
Jooken (n = 400, 1 000) & 0,0012 & 0,0010 & 0,0010 & 0,0009 &
\textbf{0,0008} & \textbf{0,0008} \\
\textbf{Moyenne} & 0,0407 & 0,0245 & 0,0294 & 0,0259 & \textbf{0,0174} &
0,0248 \\
Rang moyen (Friedman) & 5,00 & 3,19 & 3,71 & 3,11 & \textbf{2,81} &
3,18 \\
\end{longtable}
}

Friedman : χ² = 40,07, p = 1,5·10⁻⁷ (N = 31, CD = 1,35).

\begin{itemize}
\tightlist
\item
  Le \textbf{QL pré-entraîné} obtient la \textbf{meilleure moyenne} et
  le \textbf{meilleur rang} ; il bat significativement l'AG standard (p
  Holm = 0,001, 5/26/0). Face aux autres variantes adaptatives, les
  différences ne sont \textbf{pas significatives} sur ces 31 instances
  (p Holm ≥ 0,21) : avec 5 instances par classe (tailles de 100 à 2 000
  objets), la variance entre instances est forte, et les instances de
  Jooken sont résolues à moins de 0,003 \% par toutes les variantes, ce
  qui laisse peu de place aux différences. Le gain le plus net du
  pré-entraînement porte sur la classe WC (0,083 \% contre 0,116--0,157
  \% pour les autres variantes adaptatives).
\item
  \textbf{L'oracle n'est plus le meilleur} : UX + BF3, choisi sur les
  instances de calibration générées, se classe au niveau de l'aléatoire
  sur ce jeu (0,0248 \%). Une configuration fixe réglée sur une
  distribution ne se transfère pas nécessairement ; la \textbf{politique
  apprise}, qui varie les opérateurs selon l'état de la recherche, se
  transfère mieux (0,0174 \%). Ce constat rejoint le théorème « pas de
  repas gratuit » : aucun réglage fixe n'est meilleur sur toutes les
  distributions d'instances \autocite{wolpert_macready_1997}.
\item
  \textbf{QL en ligne ≈ aléatoire} (p = 1,00), confirmant H2 sur un
  second jeu.
\end{itemize}

\textbf{Le mur des méthodes exactes (instances de Jooken,
\(C = 10^8\)).} Pour les 8 instances de capacité \(10^8\), la
\textbf{programmation dynamique est infaisable} (vecteur de \(10^8\)
réels ≈ 0,8 Go, table de reconstruction bien au-delà), et la
\textbf{PLNE (HiGHS) ne termine pas} en 60 s sur 4 d'entre elles. Sur
ces 4 instances, \textbf{toutes les variantes d'AG, en ≈ 7 s, font mieux
que la PLNE en 60 s} :

{\def\LTcaptype{none} % do not increment counter
\begin{longtable}[]{@{}
  >{\raggedright\arraybackslash}p{(\linewidth - 8\tabcolsep) * \real{0.2000}}
  >{\raggedright\arraybackslash}p{(\linewidth - 8\tabcolsep) * \real{0.2000}}
  >{\raggedright\arraybackslash}p{(\linewidth - 8\tabcolsep) * \real{0.2000}}
  >{\raggedright\arraybackslash}p{(\linewidth - 8\tabcolsep) * \real{0.2000}}
  >{\raggedright\arraybackslash}p{(\linewidth - 8\tabcolsep) * \real{0.2000}}@{}}
\toprule\noalign{}
\begin{minipage}[b]{\linewidth}\raggedright
Instance (Jooken)
\end{minipage} & \begin{minipage}[b]{\linewidth}\raggedright
AG standard
\end{minipage} & \begin{minipage}[b]{\linewidth}\raggedright
QL pré-entraîné
\end{minipage} & \begin{minipage}[b]{\linewidth}\raggedright
PLNE 60 s
\end{minipage} & \begin{minipage}[b]{\linewidth}\raggedright
Glouton
\end{minipage} \\
\midrule\noalign{}
\endhead
\bottomrule\noalign{}
\endlastfoot
n = 1 000, g = 10, s = 100 & 0,0007 & \textbf{0,0001} & 0,0003 &
0,2464 \\
n = 1 000, g = 10, s = 300 & 0,0021 & \textbf{0,0020} & 0,0253 &
0,4662 \\
n = 400, g = 10, s = 100 & 0,0027 & 0,0023 & 0,0092 & 1,0400 \\
n = 400, g = 10, s = 300 & 0,0012 & \textbf{0,0011} & 0,0013 & 0,2244 \\
\end{longtable}
}

(écarts en \%, moyenne de 10 exécutions pour les AG). C'est la situation
où une métaheuristique garde tout son intérêt face aux méthodes exactes
(H4), comme le soulignent les revues récentes sur le sac à dos
\autocite{cacchiani_knapsack_2022-1,wilbaut_knapsack_2022}.

\textbf{Figures associées} : \texttt{fig02\_e2\_ecart\_par\_classe},
\texttt{fig02\_e2\_rangs\_friedman}, \texttt{fig02\_e2\_boites},
\texttt{fig02\_e2\_convergence}, \texttt{fig02\_e2\_succes\_echecs},
\texttt{fig02\_e2\_qualite\_temps}.

\section{Cas d'usage réels (données synthétiques illustratives ; 5
exécutions par variante
d'AG)}\label{analyse_resultats-cas-dusage-ruxe9els-donnuxe9es-synthuxe9tiques-illustratives-5-exuxe9cutions-par-variante-dag}

{\def\LTcaptype{none} % do not increment counter
\begin{longtable}[]{@{}
  >{\raggedright\arraybackslash}p{(\linewidth - 14\tabcolsep) * \real{0.1250}}
  >{\raggedright\arraybackslash}p{(\linewidth - 14\tabcolsep) * \real{0.1250}}
  >{\raggedright\arraybackslash}p{(\linewidth - 14\tabcolsep) * \real{0.1250}}
  >{\raggedright\arraybackslash}p{(\linewidth - 14\tabcolsep) * \real{0.1250}}
  >{\raggedright\arraybackslash}p{(\linewidth - 14\tabcolsep) * \real{0.1250}}
  >{\raggedright\arraybackslash}p{(\linewidth - 14\tabcolsep) * \real{0.1250}}
  >{\raggedright\arraybackslash}p{(\linewidth - 14\tabcolsep) * \real{0.1250}}
  >{\raggedright\arraybackslash}p{(\linewidth - 14\tabcolsep) * \real{0.1250}}@{}}
\toprule\noalign{}
\begin{minipage}[b]{\linewidth}\raggedright
Cas
\end{minipage} & \begin{minipage}[b]{\linewidth}\raggedright
Glouton
\end{minipage} & \begin{minipage}[b]{\linewidth}\raggedright
PD
\end{minipage} & \begin{minipage}[b]{\linewidth}\raggedright
PLNE
\end{minipage} & \begin{minipage}[b]{\linewidth}\raggedright
AG standard
\end{minipage} & \begin{minipage}[b]{\linewidth}\raggedright
Aléatoire
\end{minipage} & \begin{minipage}[b]{\linewidth}\raggedright
QL en ligne
\end{minipage} & \begin{minipage}[b]{\linewidth}\raggedright
QL pré-entraîné
\end{minipage} \\
\midrule\noalign{}
\endhead
\bottomrule\noalign{}
\endlastfoot
Budget communal (KP, n = 120) & 0,073 \% & non applicable (coûts
décimaux) & \textbf{optimal, 0,07 s} & 0,003 \% & \textbf{0 \%} &
\textbf{0 \%} & \textbf{0 \%} \\
Camion humanitaire (MKP, m = 2, n = 250) & 2,05 \% & non applicable &
\textbf{optimal, 0,3 s} & 0,096 \% & 0,070 \% & 0,061 \% & 0,061 \% \\
Machines virtuelles (MKP, m = 3, n = 300) & 25,4 \% & non applicable &
\textbf{optimal, 1,5 s} & \textbf{2,29 \%} & 4,05 \% & 2,84 \% & 3,36
\% \\
Portefeuille (KP corrélé, n = 400) & 0,173 \% & optimal, 0,17 s &
optimal, 1,7 s & 0,0045 \% & 0,0019 \% & 0,0022 \% & \textbf{0,0014
\%} \\
\end{longtable}
}

Ces cas reprennent des applications classiques du sac à dos : sélection
de projets sous contrainte budgétaire, chargement multi-ressources,
allocation de ressources informatiques et choix de portefeuille
\autocite{kellerer_knapsack_book_2004,vaezi_mean-var_2020,wilbaut_knapsack_2022}.

\begin{itemize}
\tightlist
\item
  Sur ces tailles modestes, la \textbf{PLNE} est la méthode de choix
  (optimale en moins de 2 s).
\item
  Le \textbf{glouton} est très fragile en multidimensionnel (25 \%
  d'écart sur les machines virtuelles), faute d'un ordre unique des
  efficacités lorsque plusieurs ressources sont contraintes
  \autocite{kellerer_multidimensional_2004}.
\item
  Le QL pré-entraîné est le meilleur des AG sur 3 cas sur 4 ; il
  \textbf{échoue} sur les \textbf{machines virtuelles}, où l'AG
  standard, moins disruptif, fait mieux (voir la discussion).
\end{itemize}

\section{E3 --- Sac à dos multidimensionnel, OR-Library (58 instances ×
6 variantes × 10
exécutions)}\label{analyse_resultats-e3-sac-uxe0-dos-multidimensionnel-or-library-58-instances-6-variantes-10-exuxe9cutions}

Ce jeu, issu de l'OR-Library
\autocite{beasley_orlib_1990,chu_beasley_1998}, teste le
\textbf{transfert inter-problèmes} : la Q-table Q\_OOD, apprise sur le
KP 0/1, est appliquée \textbf{sans modification} au MKP (\(m\) = 5 ou 10
contraintes). Référence : optimum SAC-94 (\emph{weing}) ou meilleure
valeur connue (PLNE 60 s ou meilleure exécution).

\textbf{Tableau 4 --- Écart moyen à la référence (\%)} :

{\def\LTcaptype{none} % do not increment counter
\begin{longtable}[]{@{}
  >{\raggedright\arraybackslash}p{(\linewidth - 16\tabcolsep) * \real{0.1111}}
  >{\raggedright\arraybackslash}p{(\linewidth - 16\tabcolsep) * \real{0.1111}}
  >{\raggedright\arraybackslash}p{(\linewidth - 16\tabcolsep) * \real{0.1111}}
  >{\raggedright\arraybackslash}p{(\linewidth - 16\tabcolsep) * \real{0.1111}}
  >{\raggedright\arraybackslash}p{(\linewidth - 16\tabcolsep) * \real{0.1111}}
  >{\raggedright\arraybackslash}p{(\linewidth - 16\tabcolsep) * \real{0.1111}}
  >{\raggedright\arraybackslash}p{(\linewidth - 16\tabcolsep) * \real{0.1111}}
  >{\raggedright\arraybackslash}p{(\linewidth - 16\tabcolsep) * \real{0.1111}}
  >{\raggedright\arraybackslash}p{(\linewidth - 16\tabcolsep) * \real{0.1111}}@{}}
\toprule\noalign{}
\begin{minipage}[b]{\linewidth}\raggedright
Sous-ensemble
\end{minipage} & \begin{minipage}[b]{\linewidth}\raggedright
AG standard
\end{minipage} & \begin{minipage}[b]{\linewidth}\raggedright
Aléatoire
\end{minipage} & \begin{minipage}[b]{\linewidth}\raggedright
UCB1
\end{minipage} & \begin{minipage}[b]{\linewidth}\raggedright
QL en ligne
\end{minipage} & \begin{minipage}[b]{\linewidth}\raggedright
QL pré-entraîné
\end{minipage} & \begin{minipage}[b]{\linewidth}\raggedright
Oracle fixe
\end{minipage} & \begin{minipage}[b]{\linewidth}\raggedright
PLNE 60 s
\end{minipage} & \begin{minipage}[b]{\linewidth}\raggedright
Glouton
\end{minipage} \\
\midrule\noalign{}
\endhead
\bottomrule\noalign{}
\endlastfoot
OR5x100 (30) & 0,252 & 0,164 & 0,161 & 0,173 & 0,159 & \textbf{0,127} &
0,000 & 3,56 \\
OR10x100 (10) & 0,808 & 0,578 & 0,588 & 0,606 & 0,629 & \textbf{0,526} &
0,037 & 7,43 \\
OR5x250 (10) & 0,345 & 0,229 & 0,224 & 0,230 & 0,211 & \textbf{0,184} &
0,015 & 4,39 \\
SAC-94 \emph{weing} (8) & 0,041 & \textbf{0,025} & 0,036 & 0,043 & 0,062
& \textbf{0,025} & 0,000 & 2,11 \\
\textbf{Moyenne} & 0,335 & 0,228 & 0,228 & 0,240 & 0,236 &
\textbf{0,191} & --- & --- \\
Rang moyen (Friedman) & 5,55 & 3,36 & 3,17 & 3,69 & 3,29 & \textbf{1,93}
& --- & --- \\
\end{longtable}
}

Friedman : χ² = 128,2, p = 5,7·10⁻²⁶ (N = 58, CD = 0,99). Temps moyen
par exécution d'AG : ≈ 1,9 s.

\begin{itemize}
\tightlist
\item
  \textbf{H1 confirmée sur le MKP} : toutes les variantes adaptatives
  battent l'AG standard (QL pré-entraîné : p Holm \textless{} 10⁻⁵,
  28/30/0 ; QL en ligne : 22/36/0).
\item
  \textbf{Le transfert inter-problèmes n'apporte pas de gain
  supplémentaire} : QL pré-entraîné, aléatoire, UCB1 et QL en ligne ne
  se distinguent pas (p Holm ≥ 0,41 ; seul UCB1 est très légèrement
  meilleur que le QL en ligne, p = 0,037). La politique apprise sur le
  KP 0/1 conserve la diversification mais ne capture pas la structure
  multi-contraintes, que les travaux dédiés au MKP intègrent
  explicitement dans l'état ou dans les caractéristiques de l'instance
  \autocite{bushaj_k-means_2024,rezoug_application_2022,cacchiani_knapsack_2022}.
  L'\textbf{oracle} (UX + BF3) reste significativement meilleur (0/51/7)
  : le croisement uniforme est pertinent pour le MKP, mais la politique
  transférée l'utilise moins systématiquement (en phase intermédiaire
  notamment).
\item
  \textbf{Méthodes exactes} : la PLNE (HiGHS) prouve l'optimalité sur 38
  instances et atteint sa limite de 60 s sur 20 (9 OR10x100, 10 OR5x250,
  1 OR5x100), avec un écart résiduel de 0,015 à 0,037 \%. À ces tailles,
  \textbf{la PLNE en 60 s reste nettement meilleure que l'AG en ≈ 2 s} ;
  l'AG offre un compromis qualité/temps (0,2 \% en 2 s) et le
  \textbf{glouton} est très médiocre en multidimensionnel (2 à 7 \%).
\end{itemize}

\textbf{Figures associées} : \texttt{fig03\_e3\_ecart\_par\_classe},
\texttt{fig03\_e3\_rangs\_friedman}, \texttt{fig03\_e3\_succes\_echecs},
\texttt{fig03\_e3\_qualite\_temps}.

\section{E5 --- Passage à l'échelle (classes UC et SC, n = 100 à 10 000,
3
graines)}\label{analyse_resultats-e5-passage-uxe0-luxe9chelle-classes-uc-et-sc-n-100-uxe0-10-000-3-graines}

\textbf{Tableau 5 --- Temps moyen (s) et écart à la meilleure valeur
trouvée (\%)} (budget de l'AG fixe : 200 générations × 100 individus) :

{\def\LTcaptype{none} % do not increment counter
\begin{longtable}[]{@{}
  >{\raggedright\arraybackslash}p{(\linewidth - 12\tabcolsep) * \real{0.1429}}
  >{\raggedright\arraybackslash}p{(\linewidth - 12\tabcolsep) * \real{0.1429}}
  >{\raggedright\arraybackslash}p{(\linewidth - 12\tabcolsep) * \real{0.1429}}
  >{\raggedright\arraybackslash}p{(\linewidth - 12\tabcolsep) * \real{0.1429}}
  >{\raggedright\arraybackslash}p{(\linewidth - 12\tabcolsep) * \real{0.1429}}
  >{\raggedright\arraybackslash}p{(\linewidth - 12\tabcolsep) * \real{0.1429}}
  >{\raggedright\arraybackslash}p{(\linewidth - 12\tabcolsep) * \real{0.1429}}@{}}
\toprule\noalign{}
\begin{minipage}[b]{\linewidth}\raggedright
Classe
\end{minipage} & \begin{minipage}[b]{\linewidth}\raggedright
n
\end{minipage} & \begin{minipage}[b]{\linewidth}\raggedright
PD
\end{minipage} & \begin{minipage}[b]{\linewidth}\raggedright
PLNE
\end{minipage} & \begin{minipage}[b]{\linewidth}\raggedright
Glouton
\end{minipage} & \begin{minipage}[b]{\linewidth}\raggedright
AG standard
\end{minipage} & \begin{minipage}[b]{\linewidth}\raggedright
QL pré-entraîné
\end{minipage} \\
\midrule\noalign{}
\endhead
\bottomrule\noalign{}
\endlastfoot
SC & 1 000 & 1,2 s · 0 \% & 80 s* · 0 \% & 0,01 s · 0,134 \% & 2,5 s ·
0,166 \% & 3,0 s · \textbf{0,021 \%} \\
SC & 2 000 & 5,8 s · 0 \% & 120 s* · 0 \% & 0,01 s · 0,080 \% & 5,5 s ·
0,270 \% & 5,9 s · \textbf{0,036 \%} \\
SC & 10 000 & infaisable & 120 s* · 0 \% & 0,06 s · \textbf{0,017 \%} &
15,8 s · 1,170 \% & 21,9 s · 0,365 \% \\
UC & 1 000 & 1,5 s · 0 \% & 0,5 s · 0 \% & 0,01 s · 0,011 \% & 4,1 s ·
0,005 \% & 4,5 s · \textbf{0,002 \%} \\
UC & 10 000 & infaisable & 33 s · 0 \% & 0,06 s · \textbf{0,0001 \%} &
36 s · 6,25 \% & 40 s · 1,44 \% \\
\end{longtable}
}

* limite de temps atteinte sans preuve d'optimalité. « Infaisable » :
\(n \cdot C > 3 \cdot 10^9\) opérations (garde-fou ; la PD dépasserait
plusieurs minutes et plusieurs gigaoctets pour reconstruire la
solution).

\begin{itemize}
\tightlist
\item
  La \textbf{PD} croît en \(O(nC)\), soit \(O(n^2)\) ici
  (\(C \propto n\)) : 6 s à n = 2 000, hors de portée au-delà de 5 000.
  Elle est pseudo-polynomiale : le KP 0/1 n'est NP-difficile qu'au sens
  faible \autocite{garey_johnson_1979,kellerer_improved_2004}.
\item
  La \textbf{PLNE} trouve la meilleure valeur à toutes les tailles, mais
  \textbf{ne prouve plus l'optimalité} sur la classe SC dès n = 500
  (limite de 80--120 s).
\item
  À budget fixe, \textbf{l'AG se dégrade avec la taille} ; le \textbf{QL
  pré-entraîné reste 3 à 8 fois meilleur que l'AG standard} jusqu'à n =
  10 000 --- soit 50 fois la taille des instances d'entraînement : le
  transfert en taille fonctionne.
\item
  \textbf{Échec à grande échelle} : à n ≥ 5 000, le \textbf{glouton}
  (0,06 s) fait mieux que toutes les variantes d'AG : avec 200
  générations et une population initiale aléatoire, l'AG n'a pas le
  temps de converger sur 10 000 variables, alors que la solution
  gloutonne est quasi optimale lorsque n est grand (\(C\) grand devant
  chaque poids), propriété liée à la borne de Dantzig
  \autocite{dantzig_1957,martello_toth_1990}. Remèdes évidents, non
  testés faute de temps : \textbf{initialisation gloutonne} d'une partie
  de la population (option \texttt{init\_greedy\_fraction} déjà
  implémentée) et \textbf{budget proportionnel à n}.
\end{itemize}

\textbf{Figures associées} : \texttt{fig07\_scalabilite}.

\printbibliography[title=Bibliographie]

\end{document}
